% GENERATED FILE. Edit canonical lessons, then run npm run generate:books.
% canonical-source-hash: fnv1a64:6185932e6e04853f
% canonical-lessons: UR-C317-id, UR-C317-hilal, UR-C317-mehndi, UR-C317-jagmagana, UR-C317-raunaq

\chapter{The Night Before Eid}
\label{ch:ur-the-night-before-eid}
\hlchaptermodality{317}

\begin{chapteropening}
\textbf{By the end of this chapter:} \emph{I can describe the evening before a festival.}

Complete the last lesson of chapter 317: I can describe the evening before a festival.
\end{chapteropening}

\section[\texorpdfstring{\ur{عید} (īd)}{عید (īd)}]{\ur{عید} (īd) — the festival of Eid}

\label{lesson:UR-C317-id}



\begin{quote}
Before the new one: say the Urdu for to slip away, then the Urdu for to crawl.
\end{quote}

\subsection*{You'll want to know: \ur{عید}}
\textbf{\ur{عید}} — \emph{īd} — \textquotedblleft{}the festival of Eid\textquotedblright{}.

\begin{quote}
\ur{کل} \ur{عید} \ur{ہے۔}
\end{quote}

\begin{quote}
\emph{kal īd hai.}
\end{quote}

\begin{quote}
Tomorrow is Eid.
\end{quote}

Feminine.

\begin{cousinweb}
\textbf{\ur{عید}} (\emph{īd}) is usually traced, through Persian, to Arabic \emph{ʿīd}, \textquotedblleft{}a festival, a day that comes round again\textquotedblright{}. It is feminine: \textbf{\ur{عید} \ur{آئی}}.
\end{cousinweb}

\begin{grammarlens}[title={What you've built}]
One more word for telling what happened.
\end{grammarlens}

\subsection*{Your turn}
\begin{itemize}
\raggedright
  \item \emph{Say it:} \emph{īd}
  \item \emph{Say it:} \emph{īd}, once more
  \item \emph{Say it:} \emph{kal īd hai.} — Tomorrow is Eid.
  \item \emph{Recall:} say the Urdu for to slip away, then the Urdu for to crawl, then say \emph{īd} again
\end{itemize}

\begin{tcolorbox}[breakable,colback=teal!4,colframe=teal!35!black,title={Before you move on}]
What does \ur{عید} mean? (\textquotedblleft{}The festival of Eid\textquotedblright{}.) Say it once more.
\end{tcolorbox}
\section[\texorpdfstring{\ur{ہلال} (hilāl)}{ہلال (hilāl)}]{\ur{ہلال} (hilāl) — the new crescent moon}

\label{lesson:UR-C317-hilal}



\begin{quote}
Before the new one: say the Urdu for to crawl, then the Urdu for the festival of Eid.
\end{quote}

\subsection*{You'll want to know: \ur{ہلال}}
\textbf{\ur{ہلال}} — \emph{hilāl} — \textquotedblleft{}the new crescent moon\textquotedblright{}.

\begin{quote}
\ur{شام} \ur{کو} \ur{آسمان} \ur{پر} \ur{ہلال} \ur{نظر} \ur{آیا۔}
\end{quote}

\begin{quote}
\emph{shām ko āsmān par hilāl nazar āyā.}
\end{quote}

\begin{quote}
In the evening the new moon was seen in the sky.
\end{quote}

Masculine.

\begin{cousinweb}
\textbf{\ur{ہلال}} (\emph{hilāl}) is usually traced, through Persian, to Arabic \emph{hilāl}, \textquotedblleft{}the new moon\textquotedblright{}: the thin first crescent that tells people the month has begun.
\end{cousinweb}

\begin{grammarlens}[title={What you've built}]
One more word for telling what happened.
\end{grammarlens}

\subsection*{Your turn}
\begin{itemize}
\raggedright
  \item \emph{Say it:} \emph{hilāl}
  \item \emph{Say it:} \emph{hilāl}, once more
  \item \emph{Say it:} \emph{shām ko āsmān par hilāl nazar āyā.} — In the evening the new moon was seen in the sky.
  \item \emph{Recall:} say the Urdu for to crawl, then the Urdu for the festival of Eid, then say \emph{hilāl} again
\end{itemize}

\begin{tcolorbox}[breakable,colback=teal!4,colframe=teal!35!black,title={Before you move on}]
What does \ur{ہلال} mean? (\textquotedblleft{}The new crescent moon\textquotedblright{}.) Say it once more.
\end{tcolorbox}
\section[\texorpdfstring{\ur{مہندی} (mehndī)}{مہندی (mehndī)}]{\ur{مہندی} (mehndī) — henna}

\label{lesson:UR-C317-mehndi}



\begin{quote}
Before the new one: say the Urdu for the festival of Eid, then the Urdu for the new crescent moon.
\end{quote}

\subsection*{You'll want to know: \ur{مہندی}}
\textbf{\ur{مہندی}} — \emph{mehndī} — \textquotedblleft{}henna\textquotedblright{}.

\begin{quote}
\ur{اس} \ur{کے} \ur{ہاتھوں} \ur{پر} \ur{مہندی} \ur{تھی۔}
\end{quote}

\begin{quote}
\emph{us ke hāthoṅ par mehndī thī.}
\end{quote}

\begin{quote}
She had henna on her hands.
\end{quote}

Feminine.

\begin{cousinweb}
\textbf{\ur{مہندی}} (\emph{mehndī}) is inherited: CDIAL lists the Hindi cousin \emph{mẽhdī}, the henna plant, under Sanskrit \emph{mēndhī} (CDIAL 10328). It is feminine.
\end{cousinweb}

\begin{grammarlens}[title={What you've built}]
One more word for telling what happened.
\end{grammarlens}

\subsection*{Your turn}
\begin{itemize}
\raggedright
  \item \emph{Say it:} \emph{mehndī}
  \item \emph{Say it:} \emph{mehndī}, once more
  \item \emph{Say it:} \emph{us ke hāthoṅ par mehndī thī.} — She had henna on her hands.
  \item \emph{Recall:} say the Urdu for the festival of Eid, then the Urdu for the new crescent moon, then say \emph{mehndī} again
\end{itemize}

\begin{tcolorbox}[breakable,colback=teal!4,colframe=teal!35!black,title={Before you move on}]
What does \ur{مہندی} mean? (\textquotedblleft{}Henna\textquotedblright{}.) Say it once more.
\end{tcolorbox}
\section[\texorpdfstring{\ur{جگمگانا} (jagmagānā)}{جگمگانا (jagmagānā)}]{\ur{جگمگانا} (jagmagānā) — to glitter; to be lit up}

\label{lesson:UR-C317-jagmagana}



\begin{quote}
Before the new one: say the Urdu for the new crescent moon, then the Urdu for henna.
\end{quote}

\subsection*{You'll want to know: \ur{جگمگانا}}
\textbf{\ur{جگمگانا}} — \emph{jagmagānā} — \textquotedblleft{}to glitter; to be lit up\textquotedblright{}.

\begin{quote}
\ur{بازار} \ur{روشنی} \ur{سے} \ur{جگمگا} \ur{رہا} \ur{تھا۔}
\end{quote}

\begin{quote}
\emph{bāzār raushnī se jagmagā rahā thā.}
\end{quote}

\begin{quote}
The market was glittering with lights.
\end{quote}

Past (it …): \ur{وہ} \ur{شہر} \ur{جگمگایا}, \emph{voh shahr jagmagāyā}.

\begin{cousinweb}
\textbf{\ur{جگمگانا}} (\emph{jagmagānā}) is built on \emph{jagmag}, \textquotedblleft{}glitter\textquotedblright{}, which CDIAL lists for Hindi under a form it reconstructs (CDIAL 5076). The doubled \emph{jag-mag} is the twinkle of many lights.
\end{cousinweb}

\begin{grammarlens}[title={What you've built}]
One more word for telling what happened.
\end{grammarlens}

\subsection*{Your turn}
\begin{itemize}
\raggedright
  \item \emph{Say it:} \emph{jagmagānā}
  \item \emph{Say it:} \emph{jagmagānā}, once more
  \item \emph{Say it:} \emph{bāzār raushnī se jagmagā rahā thā.} — The market was glittering with lights.
  \item \emph{Recall:} say the Urdu for the new crescent moon, then the Urdu for henna, then say \emph{jagmagānā} again
\end{itemize}

\begin{tcolorbox}[breakable,colback=teal!4,colframe=teal!35!black,title={Before you move on}]
What does \ur{جگمگانا} mean? (\textquotedblleft{}To glitter; to be lit up\textquotedblright{}.) Say it once more.
\end{tcolorbox}
\section[\texorpdfstring{\ur{رونق} (raunaq)}{رونق (raunaq)}]{\ur{رونق} (raunaq) — liveliness; festive bustle}

\label{lesson:UR-C317-raunaq}



\begin{quote}
Before the new one: say the Urdu for henna, then the Urdu for to glitter.
\end{quote}

\subsection*{You'll want to know: \ur{رونق}}
\textbf{\ur{رونق}} — \emph{raunaq} — \textquotedblleft{}liveliness; festive bustle\textquotedblright{}.

\begin{quote}
\ur{چاند} \ur{رات} \ur{کو} \ur{بازار} \ur{میں} \ur{رونق} \ur{تھی۔}
\end{quote}

\begin{quote}
\emph{chā̃d rāt ko bāzār meṅ raunaq thī.}
\end{quote}

\begin{quote}
On the night of the new moon the market was full of life.
\end{quote}

Feminine.

\begin{cousinweb}
\textbf{\ur{رونق}} (\emph{raunaq}) is usually traced, through Persian, to Arabic \emph{rawnaq}, \textquotedblleft{}brightness, beauty\textquotedblright{}; in Urdu it is the life a crowd gives a place.
\end{cousinweb}

\begin{grammarlens}[title={What you've built}]
One more word for telling what happened.
\end{grammarlens}

\subsection*{Your turn}
\begin{itemize}
\raggedright
  \item \emph{Say it:} \emph{raunaq}
  \item \emph{Say it:} \emph{raunaq}, once more
  \item \emph{Say it:} \emph{chā̃d rāt ko bāzār meṅ raunaq thī.} — On the night of the new moon the market was full of life.
  \item \emph{Recall:} say the Urdu for henna, then the Urdu for to glitter, then say \emph{raunaq} again
\end{itemize}

\begin{tcolorbox}[breakable,colback=teal!4,colframe=teal!35!black,title={Before you move on}]
What does \ur{رونق} mean? (\textquotedblleft{}Liveliness; festive bustle\textquotedblright{}.) Say it once more.
\end{tcolorbox}
